\section{Problems, networks and comparisons}\label{sec:methods}

\subsection{Problems and data}\label{sec:methods:data}

Two families of linear elastostatic problems were used, one in two dimensions and one in
three. Both are clamped on the face $x=0$ and carry four load cases that share one stiffness
matrix: a downward and an axial traction on the face opposite the clamp, a shear traction on
the top face, and a downward body force. The magnitude of each load case is drawn
independently as $0.05$ times a number uniform on $[0.5,1.5]$. Young's modulus is 1 and
Poisson's ratio is uniform on $[0.25,0.38]$ per instance; all quantities are dimensionless.
Meshes were generated with Gmsh \cite{Geuzaine2009}, and reference solutions were computed
by sparse direct factorisation.

\subsubsection{Two dimensions}
Rectangular plates $[0,l_x]\times[0,l_y]$ in plane stress, with $l_x$ uniform on $[1.5,3]$,
$l_y$ uniform on $[0.8,1.5]$ and zero to three non-overlapping circular holes of radius
between $0.06$ and $0.16$ times $\min(l_x,l_y)$, are meshed with linear triangles of target
size uniform on $[0.05,0.12]$ and assembled with scikit-fem \cite{Gustafsson2020}. The
corpus has 30,000 instances: 256 form the validation set and the first 1,024 of the
remainder the training pool. The last \numHoldTwoTail{} instances of the remainder lie
beyond the deepest prefix of it that any recorded run read, and \numHoldTwoN{} of them form
the holdout set of \cref{sec:holdout}.

\subsubsection{Three dimensions}
Boxes $[0,l_x]\times[0,l_y]\times[0,l_z]$, with $l_x$ uniform on $[1.5,3]$, $l_y$ on
$[0.8,1.5]$, $l_z$ on $[0.6,1.2]$ and zero to three non-overlapping spherical cavities of
radius between $0.08$ and $0.16$ times $\min(l_x,l_y,l_z)$, are meshed with linear
tetrahedra at a characteristic length $l_c$ drawn uniformly from $[0.0579,0.0906]$ per
instance. The corpus has 2,000 instances: 256 form the validation set, the first 1,024 of
the remainder the training pool, and the other \numHoldThreeTail{} were not used for
training or validation; \numHoldThreeN{} of them form the holdout set of \cref{sec:holdout}.
The validation meshes have between \numNodesInMin{} and \numNodesInMax{} nodes (median
\numNodesInMedian{}), with three degrees of freedom per node. A finer set of 320 instances
from the same distribution is meshed at $l_c=0.0374$, about half the element size; its first
256 instances, with between \numNodesFineMin{} and \numNodesFineMax{} nodes (median
\numNodesFineMedian{}), test transfer to a finer mesh, and the last 64 are reserved for
few-shot training.

\subsubsection{Conditioning}
On sixteen assembled validation instances of each corpus, the spectral condition number
$\kappa$ of the stiffness matrix ranged from \numKappaTwoLow{} to \numKappaTwoHigh{} in two
dimensions and from \numKappaThreeLow{} to \numKappaThreeHigh{} in three, so that
$\kappa^{1/2}$ lies between about \numSqrtKappaMin{} and \numSqrtKappaMax{}.

\subsection{Networks}\label{sec:methods:networks}

\subsubsection{Transformer}
Every mesh node is a token. Its input features are the node's coordinates, centred and
scaled to unit root-mean-square radius; its Dirichlet flags; and its consistent nodal load
for the load case, divided by $F_{\max}$, the largest absolute nodal load component over the
instance's four load cases. Two global descriptors are appended to every token: a load
summary (the total load magnitude and the net load per component, each divided by the number
of nodes and by $F_{\max}$, and the fraction of loaded nodes) and a geometry descriptor (the
extents, each divided by the upper end of its sampling range, so that the absolute size is
kept; Poisson's ratio; the number of holes or cavities and their area or volume fraction,
and in two dimensions their mean radius). An encoder of 8 pre-normalised transformer blocks
\cite{Vaswani2017} of width 256 with 8 attention heads applies full self-attention over the
nodes; a decoder maps each node's latent vector, concatenated with the mean latent vector
over the nodes, to the node's displacement through a three-layer perceptron, and the
constrained components are set to zero. Each load case is a separate sequence; the four
sequences of an instance are processed as one batch. The encoder and decoder have
\numParamsTransformer{} million parameters. The decoded field is multiplied by $F_{\max}$,
so that the prediction is proportional to the load amplitude; since $F_{\max}$, a nodal
load, shrinks with the element size, this scaling makes the output depend on the mesh
(\cref{sec:discussion}).

\subsubsection{Graph network}
As a supervised baseline of a widely used kind we used a MeshGraphNets-style
encode-process-decode network \cite{Pfaff2021} with node and edge encoders, 8
message-passing layers of width 256 over the mesh edges (edge features: relative position
and length), the same node features and the same output scaling as the transformer; it has
\numParamsGraph{} million parameters.

\subsubsection{Baselines}
The accuracy tables also report three deterministic baselines: the zero field; a
nearest-neighbour field, the solution of the labelled training instance nearest in the
geometry descriptor, interpolated to the new mesh in normalised coordinates; and a scaled
polynomial, a quadratic ridge regression of the field's direction on the node features with
its norm regressed on the instance's descriptors. The last two use the same labelled
instances as the supervised networks.

\subsection{Losses and training}\label{sec:methods:training}

The label-free transformer minimises \eqref{eq:training} with $L=4$: at each step one
training instance is drawn, without replacement within an epoch, its four load cases are
predicted, and the mean of their energies is the loss; no reference solution is read. The
displacement-supervised transformer, the same network, minimises on labelled instances the
mean over the load cases of the relative displacement error,
\begin{equation}
\mathcal{L}_D=\frac{1}{L}\sum_{j=1}^{L}\frac{\norm{u_j-\Ustar_j}_2}{\norm{\Ustar_j}_2},
\label{eq:ld}
\end{equation}
where $u_j$ is the prediction and $\Ustar_j$ the reference solution for load case $j$ of an
instance, and $F_j$ its load vector. The energy-norm-supervised transformer, again the same
network, minimises the relative energy-norm error,
\begin{equation}
\mathcal{L}_K=\frac{1}{L}\sum_{j=1}^{L}\frac{\norm{u_j-\Ustar_j}_K}{\norm{\Ustar_j}_K}
=\frac{1}{L}\sum_{j=1}^{L}g(u_j)^{1/2},
\label{eq:lk}
\end{equation}
which replaces the norm of \eqref{eq:ld} and nothing else. Since
$\norm{u_j-\Ustar_j}_K^2=u_j^{\top}Ku_j-2F_j^{\top}u_j+F_j^{\top}\Ustar_j$, the loss
$\mathcal{L}_K$ uses each reference solution only through the number
$\norm{\Ustar_j}_K^2=F_j^{\top}\Ustar_j$, and its gradient with respect to $u_j$,
\begin{equation}
\nabla_{u_j}\mathcal{L}_K=\frac{Ku_j-F_j}{L\,\norm{u_j-\Ustar_j}_K\,\norm{\Ustar_j}_K},
\label{eq:lkgrad}
\end{equation}
is that of the label-free objective of one instance, $(Ku_j-F_j)/L$, weighted per load case.
The two energy-norm losses thus differ in the weighting of the load cases and in the labels,
not in the norm. In three dimensions a fourth variant, the displacement-supervised
transformer with an energy term, adds the mean energy of the four load cases to
$\mathcal{L}_D$, with its gradient scaled to at most the norm of the gradient of
$\mathcal{L}_D$, or with a fixed weight with 16 labelled instances (supplementary
Section~\SuppPrereg); the graph network is trained on $\mathcal{L}_D$.

All networks were trained with AdamW \cite{Loshchilov2019}, with weight decay $10^{-4}$ and
the gradient norm clipped at 1, for 200 epochs. The learning rate followed a cosine schedule
after a linear warm-up over the first 5\% of the steps, to a peak of $10^{-3}$ for the
label-free and $1.5\times10^{-3}$ for the supervised networks; these rates were fixed before
the runs and not tuned on them, and the energy-norm-supervised transformer used the
displacement-supervised settings unchanged. Training was in single precision with
TensorFloat-32 (TF32) matrix products, with three seeds (0, 1, 2) per network. The networks
are implemented in PyTorch, and the report of each run records the versions of Python,
PyTorch, NumPy and SciPy with which it ran (supplementary Section~\SuppRepro{}). The
supervised networks were trained with 16, 64, 256 and 1,024 labelled instances, the first
16, 64, 256 and 1,024 instances of the training pool, the graph network with 64 and 1,024
only, a saving of machine time fixed in the pre-registrations of September and October 2026.
Every network ran 200 epochs, so that a supervised network with $n$ labelled instances took
$200n$ steps, 3,200 with 16, and the comparisons at the smaller budgets differ in the number
of steps as well as in the labels. The label-free transformer was trained on the same 1,024
instances without their solutions, for 204,800 steps; its training set is thus the
supervised networks' largest one with the labels removed. Each network is evaluated as it
stands at the end of its last epoch; no checkpoint was selected on the validation instances.

The two norms were compared with everything else held fixed. In the two-dimensional
comparison, at a given seed the transformers start from the same initial weights and, with
1,024 instances, see them in the same sequence of per-epoch permutations, so that the
displacement-supervised and the energy-norm-supervised transformers differ only in the loss.
The three-dimensional energy-norm-supervised transformer was trained in October 2026, about
seven weeks after the displacement-supervised transformer with which it is compared; at each
seed it started from that network's initial weights and saw the instances in its order, as a
regression of the two versions of the code on small problems showed. The two-dimensional
label-free transformers were trained, with these settings, in a separate pre-registered run
and reused without retraining (supplementary Section~\SuppSide).

\subsection{Metrics}\label{sec:methods:metrics}

For each validation instance and load case we compute five metrics. The relative
displacement error is $\norm{u-\Ustar}_2/\norm{\Ustar}_2$. The relative energy gap is $g(u)$
of \eqref{eq:relgap}, computed as $\gap(u)/\lvert\Pi_h(\Ustar)\rvert$; its square root
$g^{1/2}$ is the relative energy-norm error. With the element-wise von Mises stresses
$\vm(u)$ and $\vm(\Ustar)$, which are constant on linear elements, the von Mises error is
$\norm{\vm(u)-\vm(\Ustar)}_2/\norm{\vm(\Ustar)}_2$, the norms taken over the vector of
element values; the peak stress error is
$\lvert\max\vm(u)-\max\vm(\Ustar)\rvert/\max\vm(\Ustar)$; and the critical-region recall is
the share of the $n_c$ elements with the largest $\vm(\Ustar)$ that are also among the $n_c$
elements with the largest $\vm(u)$, with $n_c$ the nearest integer to a tenth of the number
of elements. Each metric is averaged over the four load cases of an instance and then over
the 256 validation instances; tables give the mean and the sample standard deviation over
the three seeds. The relative energy gap is the primary metric, as pre-registered: by
\cref{prop:stress} it is the strain energy of the stress error relative to that of the
solution. Being a squared error, it changes by the square of the factor by which the
energy-norm error changes; we quote factors of the energy-norm error where the two could be
confused. The normalisation of $g$ needs the solution, so $g$ is computed on evaluation
instances only; without the solution, the energy gives the ranking of \cref{cor:rank} and
the test of \cref{cor:zero}. The tables count as worse than the zero field the instance-seed
predictions whose relative energy gap, averaged over the four load cases, exceeds one.

\subsection{A classical approximation of the same size}\label{sec:methods:ritz}

For each instance we also computed the Ritz-Galerkin projection of the finite-element
problem onto a space of vector polynomials, from the instance's own stiffness matrix and
load vectors, without labels or training. With the instance's box
$[0,l_x]\times[0,l_y]\times[0,l_z]$ and $P_k$ the Legendre polynomial of degree $k$, each
displacement component is spanned by the functions
\begin{equation}
\phi_{abc}(x,y,z)=\frac{x}{l_x}\,P_a\Bigl(\frac{2x}{l_x}-1\Bigr)\,P_b\Bigl(\frac{2y}{l_y}-1\Bigr)\,
P_c\Bigl(\frac{2z}{l_z}-1\Bigr),\qquad a+b+c\le p-1,
\label{eq:ritzbasis}
\end{equation}
evaluated at the nodes, and in two dimensions by the same functions without the factor in
$z$; the factor $x/l_x$ makes every function vanish on the clamped face. With $\Phi$ the
matrix whose columns are these functions on the free degrees of freedom, the projection of a
load case is $u_p=\Phi\,(\Phi^{\top}K\Phi)^{+}\Phi^{\top}F$, the pseudo-inverse taken
through the eigendecomposition of $\Phi^{\top}K\Phi$ with the directions whose eigenvalue is
at most $10^{-13}$ times the largest dropped; on no instance was a direction dropped. Like
the rescaled prediction $\cstar u$, which is the projection onto the span of $u$
(\cref{prop:amp}), $u_p$ minimises $\Pi_h$ over its space. We report every degree of the
ladder that the specification fixed, $p=6$, 7, 8 and 9 in three dimensions, with
\numRitzThreeK{} unknowns, and $p=6$ and 8 in two, with \numRitzTwoK{}, and pair each
instance's projection with every network's prediction of it. This comparison was added after
the pre-registered runs and was not pre-registered. Its specification was written on 20
October 2026, after exploratory readings on instances drawn afresh from the same generators
and on the evaluation sets of a separate study, and was in the public repository before any
paired computation on the instances of this paper, which ran on 28 October 2026. The
projection needs the assembled stiffness matrix, which the surrogates do not.

\subsection{Comparisons fixed in advance}\label{sec:methods:prereg}

Weak baselines and selective reporting have led to over-optimism in machine learning for
fluid-related partial differential equations \cite{McGreivy2024}, and a fair comparison of
two neural operators needed common data and protocols \cite{Lu2022Fair}; pre-registration,
the recording of hypotheses, analyses and decision criteria in a dated document before the
data are seen, guards against these biases in other sciences \cite{Nosek2018}. Four runs
decide the claims of this paper, and each was pre-registered: the three-dimensional
comparison (22 to 28 September 2026), which trained the label-free transformer and reused
the supervised networks, the graph network among them, that its initial run had trained in
the first half of September 2026; the two-dimensional comparison (8 to 9 October 2026),
which trained the displacement-supervised and the energy-norm-supervised transformers and
the graph network; the holdout evaluation (21 October 2026), which read the claims of both
comparisons again on held-out instances without training; and the three-dimensional
energy-norm run (22 to 25 October 2026). Each was specified, before it was executed, in a
document that fixes the configuration, the metrics and the criteria, and that declares both
outcomes publishable; the documents were committed to the project's public repository before
their runs, and the code refused to start a run whose configuration did not match the
registered hash. Where a comparison of two networks over three seeds carries a verdict, it
uses a noise guard fixed in advance \eqref{eq:guard}: a difference counts only beyond the larger of 10\% and
twice its standard error over the seeds, and within the guard no difference is shown, which
is not equivalence. \Cref{tab:claims} lists the pre-registered claims with their outcomes,
and \cref{sec:holdout} the claims that the holdout evaluation, also pre-registered, read
again on held-out instances; departures from the plans that bear on a result are stated
with it, and the readings taken after the runs, and the Ritz baseline, are marked as such
where they are reported. Supplementary Section~\SuppPrereg{} gives the full record. Every
table, figure and quoted number is computed from the committed records by the repository's
generation script and checked by its tests (supplementary Section~\SuppRepro{}); that
script, like the code of the experiments, was written with the help of the generative AI
models Claude Fable 5, Claude Fable 5.1 and Claude Opus 5.5 (Anthropic). Two earlier
pre-registered runs in two dimensions, in July 2026, used an earlier version of the code,
which did not multiply the decoded field by $F_{\max}$; the main one, the July run, is
reported in supplementary Section~\SuppJuly{}, and only its test of a warm start
(\cref{sec:cost:labels}) and its transfer readings (\cref{sec:discussion}) are quoted here.
